\scanpage{593}

\mspage{2}

tel que $\Psi\supset\Phi$, on \unsure{trouve} un hom. can.

\[
k_\Phi\mathrel{\unsure{\longrightarrow|-}} k_\Psi\ ,
\]

donc si $\unsure{\mathscr F|\mathcal F}(E)=\mathscr F$ est l'ens des \deleted{\unsure{en}} filtres sur $E$, ordonnés par inclusion, on trouve un système inductif

\[
\Phi\mathrel{\unsure{\mapstochar\relbar|\longmapsto}} k_\Phi:\quad\mathscr F\longrightarrow\mathrm{Ann}
\]

On n'a pas utilisé encore que les $k_{\unsure{\xi|\zeta}}$ sont des corps, en fait il suffisait que ce soient des objets d'une catégorie $\unsure{\mathcal A|A}$ stable par \inserted{petits produits} \inserted{et par} petites $\underset{\longrightarrow}{\lim}$ filtrantes. Prenons maintenant $\mathcal A=\mathrm{Ann}$ (catégorie des anneaux unitaires, non néc. commutatifs), et supposons que les $k_\xi$ soient des corps.

\uline{Proposition 1} \inserted{Supposons que les $k_\xi$ $(\xi\in E)$ soient des \uline{corps}.$^{*}$} Pour que $k_\Phi$ soit un corps, il f et s. que $\Phi$ soit un ultrafiltre, i.e. que pour toute \deleted{partie $A$ de \deleted{\unsure{$\mathcal P(E)$}} $E$, on ait $A\in\Phi$ ou $\complement_E A$} \inserted{décomposition $E=A\amalg B$, on ait} $A\in\Phi$ ou $B\in\Phi$. (NB. Si $A\in\Phi$, on a $B\notin\Phi$, puisque sinon on aurait $A\cap B\in\Phi$ donc $\varnothing\in\Phi$, absurde.)

% slanted
\margin{$^{*}$NB. Cette condition est \unsure{nécessaire}, si on veut que $k_\Phi$ soit un corps pour \unsure{tout} ultrafiltre \unsure{p.|pr.} $\Phi$.}

Nécessité : Soit, pour \unsure{tte|tt|toute} partie $A$ de $E$,

\[
\unsure{f_A|\varphi_A}\in k_E\quad\text{défini par}\quad f_A(\xi)=
\begin{cases}
1_{k_\xi}&\text{si }\xi\in A\\
0_{k_\xi}&\text{si }\xi\in E\setminus A\ .
\end{cases}
\]

% archive sheet number: 592
